% Source: sources/returns/S6_04_MAPS_MANIFOLDS/sections/manifolds.tex:31–91
\paragraph{Theorem D.1.2 (Proved: the cusp of the cyclic cover).}
The inverse image $\widetilde W_p$ of the cusp is connected. Its
normalization has $d$ components, each isomorphic to $dP_6$. Its
integral homology is torsion free, with ranks
\[
 (b_0,b_1,b_2,b_3,b_4)(\widetilde W_p)=(1,2,d+3,2,d).
\]
Here the covering of the dual torus is the sublattice
$\mathbb Z\widehat u\oplus d\mathbb Z\widehat\gamma$ of
$\overline\Lambda=\mathbb Z\widehat\gamma\oplus\mathbb Z\widehat u$.
\begin{proof}
The presentation of the fundamental group in the supplied
\textup{[AC26, Theorem 7.17]} sends the regular-fibre lattice onto
$\mathbb Z/d$ by its primitive $\gamma$ coordinate: $12$ is a unit
modulo $d$, and $3x=\ell_1c$, $4y=\ell_2c$ express the other two
meridians as multiples of $c$. On the cusp $\pi_1(W)=\overline\Lambda$,
and the same map is surjective. Hence the inverse image is connected,
with the asserted lattice of deck periods. Since
$B_0(\widehat u)=e_1$ and $B_0(\widehat\gamma)=-e_2$, its toric
components are indexed by $\mathbb Z^2/(\mathbb Ze_1+d\mathbb Ze_2)$.

We compute the normalization resolution, retaining integral coefficients.
Let $R=\mathbb Z[t]/(t^d-1)$, where $t$ moves one component in the
$e_2$ direction. The degree-zero row of the resolution is the cellular
cochain complex of the dual torus with this $d$-sheeted marking; its
cohomology is $(\mathbb Z,\mathbb Z^2,\mathbb Z)$ in degrees $0,1,2$.
The degree-four row is $R$ in horizontal degree zero. In degree two,
write the six boundary curves on $dP_6$ in cyclic order as
\[
 E_1,\quad H-E_1-E_2,\quad E_2,\quad H-E_2-E_3,
 \quad E_3,\quad H-E_1-E_3.
\]
The restriction differential from the $d$ normalizations to the $3d$
double curves, in the basis $(H,E_1,E_2,E_3)$, is
\[
 B(t)=
 \begin{pmatrix}
 -1&-1&-1&-1\\
 1&1&1&t\\
 -t&-t&-1&-t
 \end{pmatrix}:R^4\longrightarrow R^3.
\]
Replacing the second row by its sum with the first and the third by
its difference with $t$ times the first gives rows with sole entries
$t-1$ in columns four and three. The first-row unit eliminates the
remaining first-row entries. Thus this matrix is equivalent over $R$
to a matrix with diagonal entries $1,t-1,t-1$ and one zero column.
Multiplication by $t-1$ on $R$ has kernel $\mathbb Z(1+t+\cdots+t^{d-1})$
and cokernel $\mathbb Z$. Therefore $\ker B$ is free abelian of rank
$d+2$, and $\operatorname{coker}B\simeq\mathbb Z^2$.
There can be no higher differential: all rows with odd vertical degree
are zero and the resolution has horizontal degrees only $0,1,2$.
The only extension in degree two is an extension of free abelian
groups, hence splits as an abelian group. This gives cohomology ranks
$(1,2,d+3,2,d)$ with no torsion. The universal coefficient theorem,
applied from the top degree downwards, gives the same assertion for
homology. The normalization resolution used here is exact on stalks:
at a point with $r$ local branches it is the augmented cochain complex
of the $(r-1)$-simplex, and hence computes ordinary cohomology even at
the self-identifications and the triple points.
\end{proof}
